\section{Backward antichains for unbounded counters}
\label{sec:antichain}

\subsection{Exact fragment and semantics}
Fix a finite set of control locations $Q$ and $d\ge1$ counters. A state is
$(q,x)$ with $x\in\N^d$. A transition is a tuple
\[
 e=(q,a,z,q'),\qquad a,z\in\N^d,
\]
with semantics
\begin{equation}
 (q,x)\xrightarrow{e}(q',x-a+z)
 \quad\Longleftrightarrow\quad x\ge a.
 \label{eq:netstep}
\end{equation}
All vector comparisons are componentwise. Consumption happens before
production. This is a finite-control place/transition-net representation; it
can encode a vector addition system by choosing appropriate input and output
arcs, but input and output arcs must not be collapsed without preserving the
enabling test.

A finite collection $B_q\subseteq\N^d$ specifies the unsafe target at location
$q$:
\[
 \Bad_q=\up{B_q}:=\{x\mid \exists b\in B_q,\ b\le x\}.
\]
Thus the question is whether a target marking can be \emph{covered}, not whether
a prescribed vector can be reached with exact equality in every coordinate.
For example, ``at least one process is in each of two critical sections'' is
upward-closed. A condition such as ``counter $i$ equals zero'' generally is not.

Initial states may form a finite union of families
\begin{equation}
 X_f=\left\{b_f+\sum_{j=1}^{m_f}n_jr_{f,j}\;:\;
                    n_j\in\N\right\},\qquad r_{f,j}\in\N^d,
 \label{eq:family}
\end{equation}
with a supplied control location $q_f$. Zero rays and $m_f=0$ are allowed.
All parameters are independent. Correlated or bounded parameters, negative
coefficients, and side constraints are different interfaces, not implicit
extensions of this one.

\subsection{The exact predecessor formula}
For one target threshold $b$, the set of predecessors under $e$ is itself an
upward cone. Define
\begin{equation}
 \pred_e(b)=a+(b-z)^+=\max(a,a+b-z),
 \label{eq:pred}
\end{equation}
where positive part and maximum are componentwise. The expression $a+b-z$ on
the right is interpreted over integers; $b-z$ on the left may equivalently be
natural-number truncated subtraction.

\begin{lemma}[Enabling-aware predecessor]
\label{lem:pred}
For $x,a,z,b\in\N^d$,
\[
 x\ge a\ \land\ x-a+z\ge b
 \quad\Longleftrightarrow\quad x\ge\pred_e(b).
\]
Consequently $\operatorname{Pre}_e(\up b)=\up{\pred_e(b)}$.
\end{lemma}
\begin{proof}
Work in one coordinate. Under $x\ge a$, the second inequality is equivalent to
$x\ge a+b-z$ over integers. Together the two inequalities say
$x\ge\max(a,a+b-z)=a+\max(0,b-z)$. The argument applies to every coordinate.
\end{proof}

\begin{example}[The delta-only bug]
There are two controls and one transition from control 0 to control 1, consuming
one token and producing one token. Initially the state is $(0,0)$. Every state
at control 1 is bad. The net change is zero, but the transition is disabled.
The correct backward basis is $C_1=\{0\}$ and $C_0=\{1\}$, so the initial
state is safe. A formula using only the net change would give predecessor
threshold zero and invent a nonexistent counterexample. The fixture
\code{enabling} checks exactly this case.
\end{example}

\subsection{Parameter separation without a solver}
A receipt must be disjoint from \emph{all} states in an initial family, not just
its base. This check has a simple exact form because the rays are nonnegative.

\begin{lemma}[Affine-family separation]
\label{lem:separation}
Let $X=b+\sum_j\N r_j\subseteq\N^d$, and let $c\in\N^d$.
Then $X\cap\up c=\varnothing$ if and only if there is a coordinate $i$ such that
\[
 b_i<c_i\qquad\text{and}\qquad (r_j)_i=0\text{ for every }j.
\]
\end{lemma}
\begin{proof}
Such a coordinate stays below $c_i$ for every parameter assignment, proving
separation. Conversely suppose no such coordinate exists. Write
$s_i=\sum_j(r_j)_i$. Whenever $b_i<c_i$, we have $s_i>0$. Choose
\begin{equation}
 t=\max_{i:b_i<c_i}\left\lceil\frac{c_i-b_i}{s_i}\right\rceil,
 \qquad\max\varnothing=0,
 \label{eq:commonparameter}
\end{equation}
and set every $n_j=t$. Then $b_i+t s_i\ge c_i$ in every deficient coordinate,
and nonnegative rays cannot spoil the other coordinates. This constructs a
member of $X\cap\up c$. With no rays, absence of an obstruction simply means
$b\ge c$, and the parameter tuple is empty.
\end{proof}

Different thresholds may use different obstructing coordinates. To show
$X\cap\up C=\varnothing$, the checker verifies the lemma separately for every
$c\in C$; it need not find one coordinate separating the entire union.
Likewise the common value $t$ is merely a convenient witness. It is not claimed
to minimize any initial population.

This interface accommodates an arbitrary number of waiting clients while
keeping the checker entirely within exact natural arithmetic. It does not turn
arbitrary parameterized programs into counter systems; population abstraction
still needs its own source-simulation proof.

\subsection{Search by minimal unsafe configurations}
Let $U^{(0)}=\Bad$ and repeatedly add all one-step predecessors. Each set is
upward-closed, so store only its componentwise minimal vectors at each control.
The producer keeps a queue of newly admitted vectors and an immutable
justification node for each admission.

\begin{lstlisting}
active[q] := empty antichain for every q
history := immutable nodes; pending := FIFO queue

admit(q, v, justification):
    if some w in active[q] has w <= v: return
    remove every w in active[q] with v <= w
    append (q, v, justification) to history
    put the node in active[q] and pending
    if an initial family at q covers v:
        construct its parameters and follow history to a bad node
        return the executable transition sequence

admit every target threshold, with terminal justification
while pending is not empty:
    remove a node (q', b) from pending
    if it is no longer active: continue
    for every incoming transition e : q -> q':
        admit(q, consume[e] + positive_part(b - produce[e]),
              (e, successor_node))
return the per-control active vectors
\end{lstlisting}

Removing a dominated vector does not remove its upward cone from the represented
set: the newly admitted smaller vector covers it. Skipping its queued expansion
is also sound. If $b'\le b$, then $\pred_e(b')\le\pred_e(b)$, so expansion
of the smaller active threshold eventually covers the predecessors of the
removed one. The worklist and active-map checks implement this argument rather
than merely asserting that subsumption should be safe.

\begin{theorem}[Termination and decision on the admitted fragment]
\label{thm:countercomplete}
Without operational cutoffs, the search terminates. It returns either a
backward-closed unsafe-region receipt disjoint from every initial family, or
an executable initial-family instance and finite trace reaching $\Bad$.
\end{theorem}
\begin{proof}
Order $Q\times\N^d$ by equal control and componentwise comparison. Finite
products of $\N$ are well quasi-ordered (Dickson's lemma), as is a finite
disjoint union of them~\cite{abdulla,finkel,wqo}. Suppose infinitely many
vectors were admitted. At its admission each new vector is outside the upward
closure of all earlier admitted vectors: deletion by subsumption does not
change that accumulated upward region. Hence there are no earlier and later
admissions in the same control with the earlier vector below the later one.
This contradicts the well-quasi-order property. Each admission schedules only
finitely many predecessor operations, so the worklist empties.

Initially the represented set is $\Bad$. Every admitted predecessor can reach
a previously justified upward cone, hence eventually $\Bad$, by
Lemma~\ref{lem:pred}. Conversely at termination the represented set contains
$\Bad$ and is predecessor-closed, so induction on trace length shows that it
contains every state from which $\Bad$ is reachable. It is therefore the exact
backward reachable region. Lemma~\ref{lem:separation} decides whether any
initial family intersects it and constructs an initial instance whenever it
does. The stored derivation supplies the corresponding trace.
\end{proof}

The theorem does not imply a small basis, short runtime, short counterexample,
or polynomial bit complexity. The mathematical termination argument and the
engineering resource limit have different roles. The prototype reports
\code{unknown} when its admission or predecessor budget is exhausted. A
partially saturated basis is not a safe receipt.

\subsection{A positive certificate smaller than the search}
A positive receipt supplies finite lists $C_q$; they need not be minimal,
antichains, or search-generated. The checker verifies:
\begin{align}
 &\forall b\in B_q,\quad \exists c\in C_q,\ c\le b,
 \label{eq:badincluded}\\
 &\forall e:q\to q',\ \forall b\in C_{q'},\quad
      \exists c\in C_q,\ c\le\pred_e(b),
 \label{eq:backclosed}\\
 &\forall f,\quad X_f\cap\up{C_{q_f}}=\varnothing.
 \label{eq:initseparate}
\end{align}
The final condition uses the obstruction coordinates of
Lemma~\ref{lem:separation}. The certificate can carry those coordinates as
indices for faster checking; the delivered simple checker finds them itself.

\begin{theorem}[Counter receipt soundness]
\label{thm:countersound}
Conditions~\eqref{eq:badincluded}--\eqref{eq:initseparate} imply that no finite
execution from an admitted initial state reaches $\Bad$.
\end{theorem}
\begin{proof}
Set $I(q,x):\!\iff x\notin\up{C_q}$. This holds initially. Suppose a valid step
from $(q,x)$ reaches $(q',y)$ and $y\in\up{C_{q'}}$. Choose $b\in C_{q'}$
with $b\le y$. Lemma~\ref{lem:pred} yields $\pred_e(b)\le x$;
\eqref{eq:backclosed} then gives a $c\in C_q$ below $x$, contradicting $I$.
Thus $I$ is preserved. By \eqref{eq:badincluded}, $I$ excludes every bad state.
Induct over the finite execution.
\end{proof}

This theorem needs neither Dickson's lemma nor correctness of the search.
The formalized checker can be shipped before a formal proof of search
termination. Its invariant is the \emph{complement} of a backward unsafe region,
which is often a disjunction:
\[
 I(q,x)\equiv\bigwedge_{c\in C_q}\bigvee_{i<d}x_i<c_i.
\]
For natural counters, zero thresholds contribute impossible disjuncts; no
subtraction of one from zero should be performed when presenting this formula.

If $M=\sum_q|C_q|$, the naive closure check costs
\[
 O\!\left(d\sum_{e:q\to q'}|C_q|\,|C_{q'}|\right)
\]
integer comparisons and additions. Initial checks cost
$O(d\sum_f m_f|C_{q_f}|)$ with a small extra term for the bases; precomputing
which coordinates have nonzero ray sums reduces repeated work. All figures
count arithmetic operations, not bit operations. Counter magnitudes and
certificate byte limits remain explicit resources.

\subsection{Why witness history must outlive subsumption}
An unsafe receipt contains the selected initial family, its parameter tuple,
and a list of original transition indices. The checker executes that list,
checking every source control and consumption bound, and tests the final
marking against the original bad set. It does not accept a producer's claim
that a basis vector is reachable.

A backward justification has the shape
\[
 (q,\pred_e(b))\xRightarrow{e}(q',b)\xRightarrow{\text{suffix}}\Bad,
\]
where the first arrow means that firing from the threshold produces a marking
\emph{at least} $b$, not necessarily equal to it. Monotonicity makes the same
suffix executable from any larger marking. This is the reason threshold
justifications can be composed.

The chain must use persistent node identifiers, not references to mutable
antichain entries. In the one-counter fixture with increment $+1$, initial
zero, and target 101, the active basis always has size one: $101$ is replaced
by $100$, then by $99$, and so on. The successful negative receipt nevertheless
has 101 steps and uses 102 historical admissions. Deleting dominated history
would destroy the only explanation of the answer. No shortest-counterexample
claim is made for this subsumption search.

\subsection{Worked unbounded examples}
\paragraph{A nonconvex safe region.}
Start from any marking $(1+n,0)$ or $(0,1+n)$, with $n\in\N$. A transition
consumes one token from a coordinate and returns two to that same coordinate.
The unsafe target is $(1,1)$. The single basis vector $C=\{(1,1)\}$ is
backward-closed, and each initial family has a different fixed zero coordinate.
The receipt proves that both counters are never simultaneously positive,
although either counter is unbounded. A single convex region containing both
families cannot separate the bad target: it contains $(2,0)$ and $(0,2)$ and
therefore their midpoint $(1,1)$. This distinguishes the receipt from a single
polyhedral invariant, not from every possible polynomial or disjunctive proof.

\paragraph{Mutual exclusion for arbitrary waiting populations.}
Use coordinates $(\ell,a,b,u,v)$ for a lock, active clients of two kinds, and
waiting clients. The initial family is
\[
 (1,0,0,n,m),\quad n,m\in\N.
\]
Entering the first critical section consumes one lock and one waiting token of
kind $u$, and produces an active-$a$ token. Release reverses this transfer;
kind $v$ acts analogously. The bad threshold is $(0,1,1,0,0)$. The delivered
search computes six minimal backward thresholds, checks 24 predecessor
operations, and certifies all $n,m$ simultaneously. Removing lock consumption
from the second entry rule produces a two-step trace at $(n,m)=(1,1)$.

This toy protocol also has an obvious resource invariant
$\ell+a+b=1$ in its correct version. The experiment is not evidence that other
methods cannot prove it; it demonstrates parameter extraction, original-edge
replay, and a useful complete fragment without a user-supplied invariant.
